\scanpage{19}
\begin{flushright}Tribu n° 66\end{flushright}
\begin{center}
\uline{DÉCISIONS sur l' INTÉGRATION}

(Rédactions n° 446 et 447)

\unsure{-:-:-:-}
\end{center}

\uline{La réédition d'Intégration V} (n° 446).

Le projet de réédition est adopté d'enthousiasme. Comme la Tribu a dormi ou a mis la chaudière en marche pendant certaines discussions, que Meyer et Douady ont noté les décisions d'ordre rédactionnel, et qu'ils ont commencé à les mettre à exécution, on ne trouvera pas grand chose dans ce qui suit. Meyer et Douady sont en train de faire la rédaction définitive, dont de rares exemplaires circuleront. A l'impression ensuite !

On remplace la notation $\unsure{\widetilde\mu^*}$ par $\mu^0$. On dira `` modérée '' au lieu de `` tempérée'' (déf. 2, p. 4) ; vive Pinay ! On donnera un sorite des fonctions, ensembles et mesures modérées. Les pages 12 à 19 de la rédaction 446 sont supersédées par la suite.

On dira `` scalairement essentiellement intégrable '' (resp. `` scalairement intégrable '') pour la notion du bas de la p. 20. Dans la déf. 1, p. 21, on dira `` $\mu$-adéquate '' (ça généralise la notion de l'édition actuelle ; la déf. 2 est ainsi remplacée par une remarque, qui le dira).

Certains ont été effrayés par le nombre de chapeaux emboîtés, et proposent qu'on fasse des dépliants indiquant les chapeaux et leurs domaines de validité. L'idée a été repoussée.

\uline{Les mesures sur les espaces séparés} (n° 447).

A la lecture de cette rédaction Bourbaki a versé des larmes grosses comme des bouchons de carafe, regrettant de n'avoir pas limité les chap. III et IV aux mesures sur les compacts et de n'être pas ensuite passé directement aux espaces séparés généraux. Bref, c'est la ressurection triomphante des compactologies. Mais il est trop tard pour bouleverser encore les chap. III et IV. Hélas, hélas, hélas !

Comme le contenu du n° 447 n'a pas beaucoup plus à voir avec le chap. V qu'avec les autres, il formera un chapitre séparé, qui viendra à la fin du Livre (Chap. IX) ; un ``Leitfaden'' dira de quoi il dépend.

Douady et Meyer feront la rédaction suivante, qui sera tirée et éventuellement relue en Congrès. Schwartz fournira un papier sur la mesure de Wiener-\unsure{Levy|Lovy} sur $\mathcal C([0,1])$.
\mspage{18}
